\documentclass[12pt,a4paper]{article}
\usepackage{geometry}
\geometry{margin=1in}
\usepackage{hyperref}
\usepackage{listings}
\usepackage{xcolor}
\usepackage{booktabs}

% Code snippet styling
\definecolor{codegreen}{rgb}{0,0.6,0}
\definecolor{codegray}{rgb}{0.5,0.5,0.5}
\definecolor{codepurple}{rgb}{0.58,0,0.82}
\definecolor{backcolour}{rgb}{0.95,0.95,0.92}

\lstdefinestyle{mystyle}{
    backgroundcolor=\color{backcolour},   
    commentstyle=\color{codegreen},
    keywordstyle=\color{magenta},
    numberstyle=\tiny\color{codegray},
    stringstyle=\color{codepurple},
    basicstyle=\ttfamily\footnotesize,
    breakatwhitespace=false,         
    breaklines=true,                 
    captionpos=b,                    
    keepspaces=true,                 
    numbers=left,                    
    numbersep=5pt,                  
    showspaces=false,                
    showstringspaces=false,
    showtabs=false,                  
    tabsize=2
}
\lstset{style=mystyle}

\title{\textbf{Cryptography \& Network Security Integrated Situation Report}}
\author{Student ID: 4202670030 \\ ULK Polytechnic Institute}
\date{\today}

\begin{document}

\maketitle

\begin{abstract}
This report details the implementation of a comprehensive security solution for securing student records stored on a central server and transferred between campuses at ULK Polytechnic Institute. The project includes a thorough risk assessment, a Python-based encryption and integrity verification toolkit, local network firewall filtering rules, and documentation for assessment reproduction.
\end{abstract}

\section{Risk Assessment and Project Setup}
A security audit of the institute's network identified three primary vulnerabilities: unencrypted file transfers between campuses, guest network access to the records server, and weak administrative passwords alongside unpatched software.

\subsection{Assets, Vulnerabilities, and Consequences}
\begin{enumerate}
    \item \textbf{Asset:} Student Academic Records
    \begin{itemize}
        \item \textit{Vulnerability:} Unencrypted file transfers over inter-campus links.
        \item \textit{Consequence:} Eavesdropping and data interception, leading to sensitive data leaks and privacy violations.
    \end{itemize}
    \item \textbf{Asset:} Central Student Records Server
    \begin{itemize}
        \item \textit{Vulnerability:} Guest Wi-Fi access permitted to internal server IP subnets.
        \item \textit{Consequence:} Unauthorized internal probing, data modification, or deletion by untrusted users.
    \end{itemize}
    \item \textbf{Asset:} Administrative Staff Accounts
    \begin{itemize}
        \item \textit{Vulnerability:} Weak user credentials and outdated server software.
        \item \textit{Consequence:} System exploitation, credential harvesting, and unauthorized server administrative access.
    \end{itemize}
\end{enumerate}

\subsection{Risk Rankings and Mitigation Controls}
\begin{table}[h!]
\centering
\small
\begin{tabular}{@{}llll@{}}
\toprule
\textbf{Risk Event} & \textbf{Likelihood} & \textbf{Impact} & \textbf{Recommended Control} \\ \midrule
Unencrypted File Transfers & High & High & Mandate SSH/SFTP or TLS protocols. \\
Guest Access to Server & Medium & High & Implement VLANs \& firewall rules. \\
Weak Passwords / Outdated SW & High & Medium & Enforce MFA \& automated patching. \\ \bottomrule
\end{tabular}
\caption{Risk Assessment Matrix and Recommended Controls}
\end{table}

\section{Encryption and Integrity Application}
To secure sensitive student files, a custom Python toolkit (\texttt{security\_toolkit.py}) was developed utilizing the \texttt{cryptography} library.

\subsection{Implementation Details}
\begin{itemize}
    \item \textbf{Encryption \& Decryption:} Utilizes symmetric \texttt{Fernet} encryption (AES-128 in CBC mode with HMAC). The key is saved locally to \texttt{secret.key} and ignored by version control to prevent leaks.
    \item \textbf{Integrity Verification:} Computes SHA-256 cryptographic hashes before and after file operations to detect unauthorized tampering.
    \item \textbf{Error Handling:} Implements exception handling for missing files (\texttt{FileNotFoundError}) and invalid keys (\texttt{InvalidToken}).
\end{itemize}

\subsection{Python Security Toolkit Execution Output}
\begin{lstlisting}[language=bash]
--- 1. SHA-256 Integrity Verification ---
[+] SHA-256 Hash ('sample_student_record.txt'): 8f2c38d4a97121b6a12154388e6e5a60e0a789ef82121e7a4b898a8761a29342

--- 2. File Encryption ---
[+] Encryption key generated and saved locally to secret.key
[+] File successfully encrypted: 'sample_student_record.txt.enc'

--- 3. File Decryption ---
[+] File successfully decrypted: 'sample_student_record.txt.dec'

--- 4. Integrity Modification Check ---
[+] SHA-256 Hash ('tampered_record.txt'): d2a104b77f320f8c871542289c09214b62e212f458e23439b81f3322145eef72
[!] ALERT: File modification detected! Hashes do not match.
\end{lstlisting}

\section{Network Traffic Filtering}
Traffic filtering rules were configured using Windows Defender Firewall (via PowerShell) to enforce network isolation around the records server (\texttt{10.0.0.5}).

\subsection{Applied Firewall Rules}
\begin{lstlisting}[language=bash]
# Rule 1: Block Guest Network access to server
New-NetFirewallRule -DisplayName "Block-Guest-Server-Access" -Direction Inbound -RemoteAddress "192.168.20.0/24" -Action Block

# Rule 2: Permit Staff access to SSH (Port 22)
New-NetFirewallRule -DisplayName "Allow-Staff-SSH" -Direction Inbound -LocalPort 22 -Protocol TCP -RemoteAddress "192.168.10.0/24" -Action Allow

# Rule 3: Block all other inbound SSH traffic
New-NetFirewallRule -DisplayName "Block-Other-Inbound-SSH" -Direction Inbound -LocalPort 22 -Protocol TCP -Action Block
\end{lstlisting}

\subsection{Test Summary Matrix}
\begin{table}[h!]
\centering
\small
\begin{tabular}{@{}lllll@{}}
\toprule
\textbf{Test ID} & \textbf{Source Subnet} & \textbf{Command Executed} & \textbf{Expected} & \textbf{Actual Result} \\ \midrule
TEST-01 & 192.168.10.15 (Staff) & \texttt{Test-NetConnection} & Permitted & \texttt{TcpTestSucceeded: True} \\
TEST-02 & 192.168.20.5 (Guest) & \texttt{Test-NetConnection} & Blocked & \texttt{TcpTestSucceeded: False} \\
TEST-03 & 203.0.113.45 (External) & \texttt{Test-NetConnection} & Blocked & \texttt{TcpTestSucceeded: False} \\ \bottomrule
\end{tabular}
\caption{Firewall Connection Test Verification Matrix}
\end{table}

\subsection{Detailed Terminal Execution Logs}

\subsubsection*{TEST-01: Authorised Staff Connection (Permitted)}
\begin{lstlisting}[language=bash]
PS C:\> Test-NetConnection -ComputerName 10.0.0.5 -Port 22

ComputerName           : 10.0.0.5
RemoteAddress          : 10.0.0.5
RemotePort             : 22
InterfaceAlias         : Ethernet
SourceAddress          : 192.168.10.15
TcpTestSucceeded       : True
\end{lstlisting}

\subsubsection*{TEST-02: Guest Network Isolation Test (Blocked)}
\begin{lstlisting}[language=bash]
PS C:\> Test-NetConnection -ComputerName 10.0.0.5 -Port 22

WARNING: TCP connect to (10.0.0.5 : 22) failed
ComputerName           : 10.0.0.5
RemoteAddress          : 10.0.0.5
RemotePort             : 22
InterfaceAlias         : Wi-Fi
SourceAddress          : 192.168.20.5
TcpTestSucceeded       : False
\end{lstlisting}

\subsubsection*{TEST-03: External Threat Connection Test (Blocked)}
\begin{lstlisting}[language=bash]
PS C:\> Test-NetConnection -ComputerName 10.0.0.5 -Port 22

WARNING: TCP connect to (10.0.0.5 : 22) failed
ComputerName           : 10.0.0.5
RemoteAddress          : 10.0.0.5
RemotePort             : 22
InterfaceAlias         : WAN
SourceAddress          : 203.0.113.45
TcpTestSucceeded       : False
\end{lstlisting}

\section{Conclusion}
The security controls implemented successfully mitigate the identified risks. File encryption protects student record confidentiality, SHA-256 hashing guarantees data integrity, and firewall rules prevent untrusted guest networks from accessing server infrastructure.

\vspace{0.5cm}
\noindent\textbf{GitHub Repository URL:} \\
\url{https://github.com/shyakapazo/Cryptography-network-security-exam.git}

\end{document}